\chapter{Membrane proteins}
\label{ch:membrane}

Many proteins live inside a cell membrane rather than in water. Ion channels,
transporters and receptors are examples. To simulate them, you must place the
protein in a patch of membrane with water on both sides. ProPrep's Membrane
Builder does this for you. This chapter shows the route that worked for the
potassium channel KcsA (\pdb{1BL8}), and how to build a bare membrane or a
plain water box.

\begin{glance}
\glanceitem{What it does}{Places your protein in a lipid membrane, adds water
and salt around it, and writes one packed structure file.}
\glanceitem{You need it when}{Your protein spans a membrane, or you want a bare
membrane patch, or a water box built the same way for comparison.}
\glanceitem{Before you start}{Load the structure with \key{1a}. Remove ions,
crystal waters and other non-protein parts with \key{4a} (very important; see
Section~\ref{sec:membrane-filter}).}
\glanceitem{Where to find it}{\menupath{4g}{Membrane Builder}, then \key{1}
\textsf{Build membrane-protein system}.}
\glanceitem{You get}{\file{bilayer.pdb} (the packed system). After
\key{5a}: \file{bilayer.prmtop} and \file{bilayer.rst7}.}
\end{glance}

\section{What a membrane system is}

A cell membrane is a \term{lipid bilayer}: two layers of fat-like molecules
(\term{lipids}) with their oily tails facing each other. Each of the two layers
is called a \term{leaflet}. A common lipid in simulations is \term{POPC}, a
neutral lipid found in animal cell membranes. Cholesterol (called
\texttt{CHL1} in ProPrep) is often mixed in.

A membrane protein has a band of water-repelling surface that sits inside the
oily core of the bilayer. Before lipids can be packed around it, the protein
must be turned and shifted so that this band lies in the plane of the
membrane. This step is called \term{orientation}.

ProPrep does all of this by calling a separate program,
\term{packmol-memgen}, which comes with AmberTools. It packs lipids, water and
ions around the protein. The result is one large structure file. Building the
simulation files from it is a separate step (\key{5a}, see
Chapter~\ref{ch:topology}).

\begin{behindscenes}
The Membrane Builder does not minimize or simulate anything. It only places
molecules. Some atoms are left a little too close to each other; the first
minimization in the Molecular Dynamics Manager (Chapter~\ref{ch:runningmd})
relaxes them.
\end{behindscenes}

\section{Step by step: KcsA in a POPC membrane}
\label{sec:membrane-steps}

This is the route that worked from start to finish for \pdb{1BL8}, a channel
made of four identical chains. Start ProPrep with \texttt{proprep --full-menu}
in a new project folder.

\subsection{Load the protein and remove everything that is not protein}
\label{sec:membrane-filter}

The raw \pdb{1BL8} file contains three potassium ions and one water molecule
inside the channel. The Membrane Builder cannot handle these (see
Section~\ref{sec:membrane-traps}). Remove them first.

\begin{steps}
  \item Load the structure: \key{1a}, \key{1}, \key{1}, \key{1}, type
        \typed{1BL8}, then \key{1} for the asymmetric unit. Press \key{m} to
        return to the main menu.
  \item Type \key{4a}, then \key{1} (\textsf{Filter PDB structure}).
  \item At \prompt{Detect redox-active sites before filtering?} type \key{n}.
        Press \key{Enter} at the viewer question and at
        \prompt{Do you want to proceed with filtering?}.
  \item ProPrep asks about each part of the file in turn. Keep every protein
        chain with \key{r}. Discard the ions and the water with \key{d}:
\begin{screen}
Chain A - Protein
[r] Retain entire component
[s] Select specific residues
[d] Discard entire component
Choose option (r): §ans«r»

Chain A - POTASSIUM ION
...
Choose option (r): §ans«d»
...
Chain B - Water
...
Choose option (r): §ans«d»
\end{screen}
  \item The summary should list only \prompt{Amino\_acid: 97 residues} for each
        of the four chains. Answer \key{y} to
        \prompt{Do you want to apply these filters?}. ProPrep saves
        \file{filtered_structure.pdb}.
\end{steps}

\begin{watchout}
In version 1.23.0, the \prompt{Select residues} question in the PDB Filter
(option \key{s}) does not accept a range such as \typed{1-97}. ProPrep shows
\prompt{invalid literal for int() with base 10: '1-97'} and the filter run is
abandoned. Type numbers separated by commas, or \typed{all}.
\end{watchout}

\subsection{Open the Membrane Builder}

\begin{steps}
  \item From the main menu type \key{4g}, then \key{1}. You see the Membrane
        Builder's control screen:
\begin{screen}
  Current Configuration
  Protein:       filtered_structure.pdb
  Lipids:        not set
  Solvent:       TIP3P
  Salt:          0.15 M K+/Cl-
  Orientation:   Auto (memembed)
  Box padding:   15.0 Å

  SYSTEM
    1   Protein Selection        filtered_structure.pdb
    2   Lipid Composition        Required
    3   Solvent & Ions           TIP3P, 0.15 M K+/Cl-

  GEOMETRY
    4   Box & Membrane Dims      15.0 Å pad, 17.5 Å water
    5   Protein Orientation      Auto (memembed)
    6   Specialized Geometry     Curvature, CompEL...

  PACKING
    7   PACKMOL Settings         tol=2.0, nloop=20/100

  PARAMETRIZATION
    8   Force Fields             ff14SB / lipid21 / TIP3P
    9   Custom Params & Solutes  Solutes: none

    r   Review Full Configuration
    g   Generate & Run
    b   Back

Enter choice (2):
\end{screen}
  \item Check the \prompt{Protein} line. It should name
        \file{filtered_structure.pdb}, the file you made in the last step. ProPrep picks
        up the current structure on its own (option \key{1},
        \textsf{Use workspace structure}).
\end{steps}

Each numbered line is a group of settings. You visit only the ones you want to
change; the rest keep the values shown. \prompt{Required} means you must set
it before building.

\subsection{Choose the lipids}

For a first system, use a single, common lipid. POPC is a good choice.

\begin{steps}
  \item Type \key{2} (\textsf{Lipid Composition}), then \key{3}
        (\textsf{Enter raw lipid string}).
  \item At \prompt{Lipid string:} type \typed{POPC}.
  \item At \prompt{Ratio string (or Enter to skip)} press \key{Enter}. (We
        typed \typed{4}; for a single lipid the ratio makes no difference.)
\begin{screen}
Lipid string: §ans«POPC»
Ratio string (or Enter to skip) (): §ans«4»
Set: POPC (ratio: 4)
\end{screen}
\end{steps}

The lipid names are the ones packmol-memgen knows. An unknown name is refused,
for example \prompt{Unknown lipid: 'FOO'}, and nothing is set.

\subsection{Set the box size}

The \term{box padding} is the strip of lipid kept between the protein and the
side of the box. The \term{water layer} is the thickness of water above and
below the membrane. Bigger values give a safer system but more atoms and a
longer build.

\begin{steps}
  \item Type \key{4} (\textsf{Box \& Membrane Dims}).
  \item Type \key{1}, then the padding. We used \typed{8}. The default is 15.
  \item Type \key{2}, then the water thickness. We used \typed{10}. The
        screen notes that \prompt{17.5 Å is typical}.
  \item Type \key{b} to go back.
\end{steps}

\begin{screen}
Select [1/2/3/4/5/6/7/b] (b): §ans«1»
Box padding (Å) (min: 0.0, default: 15.0) (15.0): §ans«8»
...
Select [1/2/3/4/5/6/7/b] (b): §ans«2»
Water layer thickness (Å) (min: 0.0, default: 17.5) (17.5): §ans«10»
\end{screen}

\begin{tip}
Our small values (8 and 10~\AA) kept the build short for this manual. For real
work, values between 10--15~\AA{} for the padding and 15--17.5~\AA{} for the
water are a safer start.
\end{tip}

\subsection{Choose how the protein is oriented}

ProPrep offers two automatic methods. \term{MEMEMBED} is the default and
worked well for KcsA. \term{PPM3} is the method behind the \term{OPM database}
(a public collection of membrane proteins already placed in a membrane). With
PPM3, ProPrep can also compare its result with OPM's.

\begin{steps}
  \item Type \key{5} (\textsf{Protein Orientation}).
\begin{screen}
    1   Automatic (MEMEMBED)
    2   Automatic (PPM3)              the method behind the OPM database
    3   Pre-oriented (skip alignment)
    4   MEMEMBED options
    5   N-terminus orientation
    6   Compare with the OPM database   Currently: Yes
    b   Back

Select [1/2/3/4/5/6/b] (1): §ans«1»
Set: Automatic orientation via MEMEMBED
\end{screen}
  \item Type \key{1} for MEMEMBED (or \key{2} for PPM3), then \key{b}.
\end{steps}

When we chose PPM3 for the same protein, ProPrep checked its placement against
OPM and printed this report:

\begin{screen}
  Asking the OPM database for 1BL8...
  OPM has 1BL8. Against its placement (2824 atoms paired with OPM's file):
    same way up as in OPM
    membrane normal: 0.1° away from OPM's
    protein centre along the normal: z = +5.2 Å here, +4.9 Å in OPM
    (difference +0.3 Å)
    OPM's hydrophobic thickness: 32.4 Å
\end{screen}

Small differences like these mean the orientation agrees with OPM. The check
needs an internet connection and only reports; it changes nothing.

\subsection{Choose the force fields}
\label{sec:membrane-ff}

A \term{force field} is the set of rules and numbers that describes how atoms
push and pull on each other (Chapter~\ref{ch:forcefields}). Here you choose
one for each kind of molecule. ProPrep passes your choices on to the Topology
Generator, which will not ask again.

\begin{steps}
  \item Type \key{8}, then \key{1} (\textsf{Select all force fields (guided
        walkthrough)}).
  \item Answer each question. For KcsA we chose:
\begin{screen}
Select protein forcefield (2): §ans«3»
Select modified amino acids ():
Select dna forcefield (2): §ans«1»
Select rna forcefield (2): §ans«1»
Select carbohydrate forcefield (2): §ans«1»
Select lipid forcefield (2): §ans«2»
Select small molecule forcefield (2): §ans«1»
Select water model (6): §ans«6»
Select divalent+ ion parameters (6): §ans«1»

Selected Forcefields:
  protein: ff14SB
  ...
  lipids: lipid21
  ...
  water: TIP3P
\end{screen}
\end{steps}

In plain words: \key{3} is ff14SB for the protein, \key{1} means ``None'' for
DNA, RNA, sugars and small molecules, \key{2} is lipid21 for the lipids, and
\key{6} is the TIP3P water model that matches ff14SB. The divalent-ion
question was answered \key{1} (None) because KcsA has no metal ions after
filtering. The menu marks the recommended choice with \prompt{(Recommended)};
the protein default is ff19SB, which goes with OPC water.

\subsection{Review and build}

\begin{steps}
  \item Type \key{r} to read the full configuration, or \key{g} to build. Both
        show the review panel; \key{g} then asks for confirmation.
  \item The panel ends with the \prompt{Equivalent command}: the exact
        packmol-memgen command ProPrep will run.
\begin{screen}
│ MEMBRANE
│   Lipids (symmetric): POPC
│   Ratio:              4
│ BOX
│   XY padding:         8.0 Å
│   Water layer:        10.0 Å
...
│ Equivalent command:
│   packmol-memgen -p filtered_structure.pdb --notprotonate --nottrim
│   -l POPC -r 4 --dist 8.0 --dist_wat 10.0 --salt --salt_override
│   --ffprot ff14SB --fflip lipid21 --ffwat tip3p -o bilayer.pdb ...

Proceed with build? [y/n] (y): §ans«y»
\end{screen}
  \item Type \key{y}. ProPrep first adds hydrogen atoms to the protein, then
        orients it, then packs the lipids, water and ions.
\end{steps}

The packing step prints its progress in plain English, one line every few
\term{loops} (rounds of rearranging molecules):

\begin{screen}
  Packing: All components together, up to 100 loops
    packmol-memgen does not wait for a perfect packing: ... "lowest so far"
    is the number to watch. What overlap is left, the minimization relaxes.
    loop 83 of 100 | objective 6.87e+03 (lowest so far 6.39e+03) | ...
    | 4.2 s per loop | at most 1 min 11 s more if every loop is needed
\end{screen}

The \term{objective} measures how badly molecules still overlap. It goes up
and down; this is normal. At the end you see:

\begin{screen}
Build successful!
  Output PDB: bilayer.pdb
  Box: 69.9 × 69.9 × 78.5 Å
  Water molecules: 5,142
  Ions: 9 K+, 13 Cl-

Workspace updated. Next steps:
  1. Run Topology Generator to parametrize the system
  2. Run Molecular Dynamics Manager for minimization/equilibration
\end{screen}

A message that packmol \prompt{used all its loops and stopped short of a
perfect packing} is also normal. The remaining close contacts are removed by
the first minimization.

\subsection{Make the simulation files}

\begin{steps}
  \item Type \key{m}, then \key{5a}, then \key{2} (\textsf{Generate tLEaP
        input for single state}).
  \item ProPrep reuses your Membrane Builder choices and does not ask about
        force fields or water:
\begin{screen}
Membrane system — using forcefields selected in membrane builder:
  protein: ff14SB
  lipids: lipid21
  water: TIP3P
Solvation skipped — system already solvated by membrane builder
\end{screen}
  \item At \prompt{Edit the tLEaP template?} type \key{n}.
  \item At \prompt{Do you want to continue without bond definitions?} type
        \key{y}. KcsA has no metal sites or disulfide bonds to define.
  \item Press \key{Enter} to accept the file name \file{bilayer_tleap.in}, and
        type \key{n} at \prompt{View the tLEaP input file?}.
  \item Type \key{4} (\textsf{Generate prmtop/rst7 files}), then \key{1}
        (\textsf{Generate missing only}).
\end{steps}

tLEaP finished with \prompt{Errors = 0; Warnings = 820}. Most of the warnings
were ``close contact'' and ``name change'' messages, which are expected for a
freshly packed membrane. You now have the two files the Molecular Dynamics
Manager needs: \file{bilayer.prmtop} and \file{bilayer.rst7}
(Chapter~\ref{ch:runningmd}).

\begin{watchout}
Right after the build, the \key{5a} menu may say that topology files have
already been generated, before you have made any. It is counting a helper
file from the hydrogen-adding step (\file{pretleap_hydrogen.prmtop}). Run
option \key{4} anyway; check that \file{bilayer.prmtop} exists afterwards.
\end{watchout}

\subsection{How long it took}

\begin{tabularx}{\linewidth}{@{}lL@{}}
\toprule
\textbf{Stage} & \textbf{Time on our laptop (KcsA, 4 chains, POPC)} \\
\midrule
Hydrogens and orientation (MEMEMBED) & about 1 minute \\
Packing (\textsf{g}) & 8.5 minutes; 4--5.5 s per ``all together'' loop \\
tLEaP and checks (\textsf{5a}) & about 1.5 minutes \\
Total & about 11 minutes; 37,282 atoms in \file{bilayer.pdb} \\
\bottomrule
\end{tabularx}

\section{Other choices}

\subsection{A membrane with no protein}

You can build a bare membrane, for example to study the lipids themselves. No
structure needs to be loaded. We built a 40\,\AA{} patch of POPC with
cholesterol in a 3:1 ratio.

\begin{steps}
  \item \key{4g}, \key{1}, then \key{1} (\textsf{Protein Selection}) and
        \key{3} (\textsf{No protein (empty bilayer)}).
  \item A bare membrane needs a fixed size. Type the side length:
\begin{screen}
No protein — building empty bilayer.
Empty bilayers need an explicit x/y patch size (packmol-memgen --distxy_fix).
Membrane patch size (Å, x and y) (70.0) (70.0): §ans«40»
\end{screen}
  \item \key{2}, \key{1} (\textsf{Build custom composition}), \key{1}
        (\textsf{Symmetric}). Then add each lipid and its \term{molar ratio}
        (how many of this lipid for each of the others):
\begin{screen}
  Lipid name (or 's' to search, 'd' when done): §ans«POPC»
  1-palmitoyl-2-oleoyl-sn-glycero-3-phosphocholine (charge: 0)
  Molar ratio for POPC (min: 1, default: 1) (1): §ans«3»
...
  Lipid name (or 's' to search, 'd' when done): §ans«chl1»
  cholesterol (charge: 0)
  Molar ratio for CHL1 (min: 1, default: 1) (1): §ans«1»
...
  Lipid name (or 's' to search, 'd' when done): §ans«d»
\end{screen}
  \item In the force-field walkthrough (\key{8}, \key{1}) answer \key{1}
        (None) for the protein. Then build with \key{g}.
\end{steps}

This took about 2 minutes and gave 1,280 water molecules in a box of
40\,$\times$\,40\,$\times$\,70\,\AA. ProPrep warned that the sterol
concentration was high for so small a patch; a larger patch avoids this
warning. In \key{5a}, ProPrep says that an empty bilayer needs no bond
definitions, and builds the topology in about 40 seconds.

\subsection{A water box with no lipids}

Option \key{4} in \textsf{Lipid Composition} builds the same kind of system
without a membrane. This is useful as a matched reference for comparison. We
tried it with ubiquitin (\pdb{1UBQ}).

\begin{screen}
Select (1): §ans«4»
╭───────────────────── Solvate Only (no membrane) ─────────────────────╮
│ Builds a solvated system with no lipid bilayer. packmol-memgen packs │
│ only water and ions around the solute, using the same engine and box │
│ conventions as a membrane build. Useful for preparing a matched      │
│ non-membrane reference of the same system.                           │
╰──────────────────────────────────────────────────────────────────────╯
Build a solvent-only system (no lipids)? [y/n] (y): §ans«y»
Use a cubic box? [y/n] (n): §ans«n»
Solvate-only mode enabled (no lipids).
\end{screen}

With a 10\,\AA{} padding the build took about 1 minute and gave 7,884 water
molecules and 21 K$^+$ / 21 Cl$^-$ ions. The output file is still called
\file{bilayer.pdb}, and the review still lists \prompt{leaprc.lipid21}, even
though there are no lipids. You can ignore both.

\subsection{Mixed and asymmetric membranes}

Under \key{2} you can also choose \textsf{Asymmetric} (different lipids in the
upper and lower leaflet), or type a raw string such as \typed{POPC:CHL1} (a
mixture) or \typed{POPC//POPE} (upper//lower). Option \key{2}
(\textsf{Search / browse lipid library}) lists the 4055 lipids packmol-memgen
knows, by category or by name (searching \typed{cholesterol} finds
\texttt{CHL1}).

\subsection{Options at a glance}

\begin{tabularx}{\linewidth}{@{}lL@{}}
\toprule
\textbf{Setting} & \textbf{What you can change} \\
\midrule
\key{1} Protein Selection & Loaded structure, another PDB file, or no protein. \\
\key{2} Lipid Composition & Custom mix, library search, raw string, or no lipids. \\
\key{3} Solvent \& Ions & Water model, salt on/off, salt concentration (default
  0.15\,M), ion types, neutralization. \\
\key{4} Box \& Membrane Dims & Padding, water layer, leaflet width, fixed
  sizes. \\
\key{5} Protein Orientation & MEMEMBED, PPM3, pre-oriented (for a file
  already placed, such as an OPM download), which way the N-terminus points,
  OPM cross-check. \\
\key{6} Specialized Geometry & Curved membranes, channel plug, double
  membrane with a charge imbalance, and other advanced options. \\
\key{7} PACKMOL Settings & Packing tolerance, number of loops,
  troubleshooting options, and a time limit for unattended builds. \\
\key{8} Force Fields & One guided walkthrough for all molecule types. \\
\key{9} Custom Params \& Solutes & Extra molecules to add to the box. \\
\bottomrule
\end{tabularx}

\medskip
Options \key{6}, \key{7} and \key{9} were opened while testing, but no build
was made with them. Leave them at their defaults unless you know packmol-memgen.

\section{If something goes wrong}
\label{sec:membrane-traps}

\subsection{The build runs for minutes and is then thrown away}

Our first KcsA build used the raw \pdb{1BL8} file, with its potassium ions and
water. ProPrep could not add hydrogens to it:

\begin{screen}
tLEaP could not add the hydrogens. packmol-memgen can add them itself
(pdb2pqr/reduce), but it then removes residues it does not recognise
(cofactors such as hemes, metal sites) and assigns protonation states of its
own, replacing the ones in this structure.
Build anyway with packmol-memgen's own protonation? [y/n] (n): §ans«y»
\end{screen}

We answered \key{y}. The orientation and packing then ran for about
7~minutes, and at the very end ProPrep refused the result:

\begin{screen}
The protein in bilayer.pdb is not the one packmol-memgen was given: atom 6 is
ALA HA in residue 1 of the file, where the protein given has LEU N in
residue 2. The Topology Generator's bond directives would not apply to it.
The build is not recorded.
\end{screen}

\begin{watchout}
If you see \prompt{tLEaP could not add the hydrogens}, answer \key{n}. Go
back, remove ions, waters and other non-protein parts with \key{4a}
(Section~\ref{sec:membrane-filter}), and build again. A \file{bilayer.pdb} left
on disk by a refused build is not used by ProPrep.
\end{watchout}

\subsection{Your settings have disappeared}

\begin{watchout}
In version 1.23.0, the Membrane Builder forgets your lipid, box and
orientation choices, without telling you, when the structure changes. This
happened after filtering with \key{4a}, and also when we opened \key{4g} again
after a finished build. Before you press \key{g}, always read the
\prompt{Current Configuration} block or the review (\key{r}).
\end{watchout}

\subsection{Going back with \texttt{undo}}

You can type \typed{undo} at any question in the Membrane Builder, as
everywhere in ProPrep (Chapter~\ref{ch:sessions}). We did this inside the
force-field walkthrough. ProPrep then replayed the whole session from the
start, which took only seconds but filled several screens with
\prompt{[REPLAY: ...]} lines. Answers after the rewind point are set aside.

\begin{watchout}
After an \texttt{undo}, check the review (\key{r}) before you build. While
testing for this manual, one earlier setting (a PACKMOL time limit) was later
found back at \prompt{none}. We could not tell for certain whether the
\texttt{undo} or the reset described above caused it.
\end{watchout}

\subsection{Other messages}

\begin{itemize}
  \item \prompt{Please enter a valid number} or \prompt{Value must be at least
        0.0}: you typed text or a negative value. Type it again.
  \item \prompt{Experimental value for POPC volume not available}: a normal
        note from packmol-memgen. Ignore it.
  \item Hundreds of tLEaP warnings in \key{5a}: normal for a packed membrane,
        as long as the summary says \prompt{Errors = 0}.
  \item In the screens you may see \texttt{packmol-memgen} and
        \texttt{teLeap} taken from a different conda environment. ProPrep
        uses the programs in the folder named by \texttt{AMBERHOME} first. If
        a build fails oddly, check which AmberTools that is
        (Chapter~\ref{ch:troubleshooting}).
\end{itemize}
